% --- BEGIN PCA + LDA PRE-PROCESSING ANALYSIS ---

\section{PCA + LDA analysis}

\subsection*{Apply PCA as a pre-processing step for LDA. Evaluate the error rate for different values of PCA dimensions ($m=1, \dots, 6$). What do you observe?}

To evaluate the impact of PCA as a pre-processing step, we first project the data onto an $m$-dimensional PCA subspace, and subsequently apply LDA to reduce it to a single dimension for classification. The default threshold (average of the projected means) is used for evaluation on the validation set.

As shown in Table \ref{tab:pca_lda_metrics}, varying the number of PCA dimensions ($m$) has virtually no impact on the classification performance. When we keep all dimensions ($m=6$), the error rate perfectly matches the baseline pure LDA model at 9.30\%. Aggressively reducing the dimensions down to $m=2, 3,$ or $4$ yields a negligible shift to 9.25\%, and reducing it to a single dimension ($m=1$) yields 9.35\%. Overall, the results are roughly identical across the board.

\begin{table}[H]
    \centering
    \begin{tabular}{|c|c|}
        \hline
        \textbf{PCA Dimensions ($m$)} & \textbf{Validation Error Rate (\%)} \\
        \hline
        1 & \textit{9.35} \\
        2 & \textit{9.25} \\
        3 & \textit{9.25} \\
        4 & \textit{9.25} \\
        5 & \textit{9.30} \\
        6 & 9.30 \\ % Matches baseline pure LDA
        \hline
    \end{tabular}
    \caption{Classification error rates using PCA as a pre-processing step for LDA across different subspace dimensions.}
    \label{tab:pca_lda_metrics}
\end{table}

\subsection*{Does PCA pre-processing improve the classification accuracy? Why or why not?}

No, PCA pre-processing does not significantly improve or degrade the classification accuracy for this dataset; we obtain roughly the same results regardless of the chosen dimensionality. 

To deeply understand why PCA does not alter the LDA classification outcome, we must analyze the mathematical relationship between their optimization objectives.

Linear Discriminant Analysis (LDA) is a supervised method that seeks to find a projection that maximizes class separability. Mathematically, it does this by maximizing the between-class scatter matrix ($S_B$) while minimizing the within-class scatter matrix ($S_W$). Conceptually, it optimizes a ratio proportional to:
$$ \text{LDA Objective} \propto \frac{S_B}{S_W} $$

Principal Component Analysis (PCA), on the other hand, is an unsupervised method that seeks to maximize the total variance of the entire dataset, represented by the total scatter matrix ($S_T$). 

The fundamental link between these two algorithms is the scatter matrix identity:
$$ S_T = S_B + S_W $$
This equation states that the total variance in the dataset ($S_T$) is exactly equal to the sum of the variance between the classes ($S_B$) and the variance within the classes ($S_W$).

If the within-class variance ($S_W$) is roughly fixed, isotropic (uniform in all directions), or relatively small, then the total scatter $S_T$ becomes heavily dominated by $S_B$. In this mathematical scenario, when PCA searches for the directions of maximum total variance ($S_T$), it is inadvertently finding the directions of maximum between-class variance ($S_B$). 

Because PCA and LDA are essentially aligning with the same fundamental variations in the data, dropping the "lowest variance" PCA dimensions only discards background noise, rather than discriminative information. When LDA is subsequently applied to this PCA-reduced subspace, it has access to the exact same discriminant directions it would have found in the raw feature space. Consequently, the pre-processing step is mathematically redundant, and the combination of PCA and LDA yields virtually identical error rates as using LDA alone.

% --- END PCA + LDA PRE-PROCESSING ANALYSIS ---